%! TeX program = lualatex
\documentclass[12pt,a4paper,oneside]{report}

\usepackage{amsmath}
\usepackage{mathtools}
\usepackage{amssymb}
\usepackage{amsfonts}
\usepackage{amsthm}
\usepackage{bm}
\usepackage{bbm}
\usepackage{centernot}
\usepackage{enumitem}

\usepackage{fontspec}
\usepackage[a4paper,margin=1.2in]{geometry}
\usepackage[colorlinks=true, allcolors=blue, plainpages=false, pdfpagelabels]{hyperref}
\usepackage{mathrsfs}

\usepackage[backend=biber,style=alphabetic,maxbibnames=99]{biblatex}
\usepackage{stmaryrd}
\SetSymbolFont{stmry}{bold}{U}{stmry}{m}{n}
\usepackage{import}
\usepackage{xifthen}
\usepackage{pdfpages}
\usepackage{booktabs}

\usepackage{graphicx}
\usepackage{tikz}
\usetikzlibrary{calc}
\usetikzlibrary{arrows.meta}
\usetikzlibrary{shapes.geometric}
\usetikzlibrary{positioning}

\usepackage{microtype}

\newcommand{\incfig}[1]{%
    \def\svgwidth{0.5\columnwidth}
    \import{./figures/}{#1.pdf_tex}
}

\makeatletter
\newtheorem*{rep@theorem}{\rep@title}
\newcommand{\newreptheorem}[2]{%
\newenvironment{rep#1}[1]{%
 \def\rep@title{#2 ##1}%
 \begin{rep@theorem}}%
 {\end{rep@theorem}}}
\makeatother

\theoremstyle{plain}
\newtheorem{theorem}{Theorem}[chapter]

\newtheorem{corollary}[theorem]{Corollary}
\newtheorem{claim}[theorem]{Claim}
\newtheorem{lemma}[theorem]{Lemma}
\newtheorem{conjecture}[theorem]{Conjecture}
\newtheorem{problem}[theorem]{Problem}
\newtheorem{proposition}[theorem]{Proposition}
\newtheorem{question}[theorem]{Question}

\newtheorem*{question*}{Question}

\theoremstyle{definition}
\newtheorem{definition}[theorem]{Definition}
\newtheorem{construction}[theorem]{Construction}
\newtheorem{remark}[theorem]{Remark}

% paired solution
\newenvironment{sol}
  {\par\noindent\textbf{Solution}.}
  {\par\bigskip}

\input{preamble/mathdefs}

\hypersetup{
  pdftitle={A Recursive Triangle Construction for Asymmetric StochTSP},
  pdfauthor={}
}

\begin{document}
\pagenumbering{roman}

\begin{titlepage}
\centering
\vspace*{0.18\textheight}
{\LARGE\bfseries A Recursive Triangle Construction\\[0.35em]
for Asymmetric \StochTSP{}\par}
\vspace{1.5em}
{\large A standalone technical note\par}
\vfill
\end{titlepage}

\clearpage
\pagenumbering{arabic}
\chapter{Recursive-Triangle Clairvoyance Gap}

This note records a natural recursive generalization of the four-vertex
directed-triangle example.  Its clairvoyance-gap ratio tends to~$4/3$.
The construction is not included in the main thesis because the layered-poset
construction gives the stronger ratio~$3/2$, while the recursive proof
requires substantially more decomposition machinery.

At depth one, take the root~$r=a$, two further permanent vertices~$b,c$, and
one independently active midpoint~$d$ with activation probability~$1/2$.
All generating arcs have unit length:
\[
  a\leftrightarrow b,\qquad
  a\leftrightarrow c,\qquad
  b\to c\to d\to b.
\]
Thus the arc $c\to b$ of a bidirected triangle has been subdivided into
$c\to d\to b$.  When $d$ is inactive the posterior tour costs~$3$, and when
it is active the posterior tour costs~$4$, so $\optpost=7/2$.  A
one-step Bellman calculation over the first probe gives $\optadapt=4$.
The depth-one gap is therefore~$8/7$.

\section{The construction}\label{sec:recursive-construction}

This section defines the recursive family of asymmetric instances.  Each block
contains two copies of the preceding block and one new stochastic midpoint,
joined by a directed three-cycle whose scale doubles with the level.  Complete
port bundles make the construction independent of the entry and exit ports
chosen by a tour or policy.

For a copy~$C$ of a recursive block, write $\mathcal P(C)$ for its
\emph{permanent vertices}, which will also serve as its \emph{ports}, and
write $\mathcal M(C)$ for its midpoint vertices.  For the canonical block
$B_k$, abbreviate these sets by $\mathcal P_k$ and $\mathcal M_k$.

The block $B_0$ consists of one permanent vertex and has no midpoint.  For
$k\ge1$, let
\[
    w_k:=2^{k-1}.
\]
Construct $B_k$ from disjoint copies $B_{k-1}^{\mathrm L}$ and
$B_{k-1}^{\mathrm R}$ and a new midpoint~$m_k$.  Add the following arcs,
all of length~$w_k$:
\begin{itemize}
    \item $x\to y$ for every
    $x\in\mathcal P(B_{k-1}^{\mathrm L})$ and
    $y\in\mathcal P(B_{k-1}^{\mathrm R})$;
    \item $y\to m_k$ for every
    $y\in\mathcal P(B_{k-1}^{\mathrm R})$; and
    \item $m_k\to x$ for every
    $x\in\mathcal P(B_{k-1}^{\mathrm L})$.
\end{itemize}
Thus, after contracting the two children, the top level is the directed
triangle
\[
    B_{k-1}^{\mathrm L}
    \longrightarrow
    B_{k-1}^{\mathrm R}
    \longrightarrow
    m_k
    \longrightarrow
    B_{k-1}^{\mathrm L}.
\]
Set
\[
    \mathcal P_k
    :=
    \mathcal P(B_{k-1}^{\mathrm L})
    \cup
    \mathcal P(B_{k-1}^{\mathrm R}),
    \qquad
    \mathcal M_k
    :=
    \mathcal M(B_{k-1}^{\mathrm L})
    \cup
    \mathcal M(B_{k-1}^{\mathrm R})
    \cup\{m_k\}.
\]
Every permanent vertex is active with probability~$1$, and every midpoint is
active independently with probability~$1/2$.

For $l\ge1$, the full recursive-block instance
$\mathcal{I}_l^{\mathrm{REC}}$ is obtained from $B_l$ by adding the
root~$r$ and, for every $x\in\mathcal P_l$, the two arcs $r\to x$ and
$x\to r$, each of length~$w_l$.  Distances in
$\mathcal{I}_l^{\mathrm{REC}}$ are the directed shortest-path distances in
this generating graph.

The generating graph of every block is strongly connected.  This follows by
induction: each child is strongly connected, and the three top-level arc
bundles connect the two children and the new midpoint in a directed cycle.
The root arcs preserve strong connectivity.  Hence all shortest-path
distances are finite and define an asymmetric metric.  Notice also that every arc
crossing the boundary of a child block is incident with one of its ports.

The block $B_l$ has $|\mathcal P_l|=2^l$ permanent vertices and
$|\mathcal M_l|=2^l-1$ midpoints.  For $1\le k\le l$, it contains
$2^{l-k}$ midpoint vertices introduced at level~$k$, each associated with
the length scale~$w_k$.

\begin{figure}[!htbp]
    \centering
    \begin{tikzpicture}[
        >=Stealth,
        block/.style={
            draw,
            dashed,
            rounded corners,
            minimum width=3.0cm,
            minimum height=1.45cm,
            align=center
        },
        every edge/.style={draw,->,thick}
    ]
        \node[block] (L) at (0,0)
            {$B_{k-1}^{\mathrm L}$\\[-0.1em]\scriptsize ports
            $\mathcal P(B_{k-1}^{\mathrm L})$};
        \node[block] (R) at (5.2,0)
            {$B_{k-1}^{\mathrm R}$\\[-0.1em]\scriptsize ports
            $\mathcal P(B_{k-1}^{\mathrm R})$};
        \node[draw,circle,minimum size=8mm] (M) at (2.6,-2.0) {$m_k$};
        \node[draw,circle,minimum size=8mm] (D) at (2.6,2.0) {$r$};

        \draw[->,thick] (L) -- node[above] {$w_k$} (R);
        \draw[->,thick] (R.south) to[bend left=12]
            node[right] {$w_k$} (M.east);
        \draw[->,thick] (M.west) to[bend left=12]
            node[left] {$w_k$} (L.south);
        \draw[->] (D) to[bend right=12] node[left] {$w_l$} (L.north);
        \draw[->] (L.north) to[bend right=12] (D);
        \draw[->] (D) to[bend left=12] node[right] {$w_l$} (R.north);
        \draw[->] (R.north) to[bend left=12] (D);
    \end{tikzpicture}
    \caption{A recursive block and the root arcs in the full instance.
    An arrow incident with a child represents the complete bundle of arcs
    incident with its ports.  The root appears only at the top level, where
    $k=l$.}
    \label{fig:recursive-asymmetric-block}
\end{figure}

\section{Statement}\label{sec:recursive-statement}

The recursive construction is designed so that both benchmarks admit exact
cost recurrences.  We first state the resulting ratio and the two cost
evaluations from which it follows.  The remaining sections then develop the
decomposition tools needed to prove those evaluations.

\begin{theorem}\label{thm:recursive-gap}
For every $l\ge1$, the instance $\mathcal{I}_l^{\mathrm{REC}}$ satisfies
\[
    \frac{\optadapt(\mathcal{I}_l^{\mathrm{REC}})}
         {\optpost(\mathcal{I}_l^{\mathrm{REC}})}
    =
    \frac{4(l+1)}{3l+4}.
\]
As $l\to\infty$, these ratios approach~$4/3$, so the asymmetric
clairvoyance gap is at least~$4/3$.
\end{theorem}

The theorem follows from the next two exact cost evaluations.

\begin{lemma}[A posteriori cost]\label{lem:post-cost}
For every $l\ge1$,
\[
    \optpost(\mathcal{I}_l^{\mathrm{REC}})=(3l+4)2^{l-2}.
\]
\end{lemma}

\begin{lemma}[Adaptive cost]\label{lem:adapt-cost}
For every $l\ge1$,
\[
    \optadapt(\mathcal{I}_l^{\mathrm{REC}})=(l+1)2^l.
\]
\end{lemma}

\begin{proof}[Proof of \cref{thm:recursive-gap}]
\Cref{lem:post-cost,lem:adapt-cost} give
\[
    \frac{\optadapt(\mathcal{I}_l^{\mathrm{REC}})}
         {\optpost(\mathcal{I}_l^{\mathrm{REC}})}
    =
    \frac{(l+1)2^l}{(3l+4)2^{l-2}}
    =
    \frac{4(l+1)}{3l+4}.
\]
The right-hand side tends to~$4/3$.
\end{proof}

\section{Decomposition tools}\label{sec:recursive-decomposition}

If the two children of every block were served in two uninterrupted visits,
the desired recurrences would follow immediately.  A general tour or policy
may instead leave a child, do work elsewhere, and later return.  We therefore
define local open-service values and show how to charge every additional
piece of a child service to an arc at its parent's level.

\subsection{Open services and expanded executions}
\label{sec:recursive-open-services}

To analyze repeated entries into a recursive child, we replace a closed tour
temporarily by an open service with freely chosen endpoint ports.  Expanding
metric movements back into generating arcs then exposes the child pieces and
top-level crossings that will be charged separately.

For a copy~$C$ of $B_k$, let $d_C$ be the directed shortest-path metric of
the generating subgraph induced by~$C$.  For a fixed active-midpoint set
$A\subseteq\mathcal M(C)$, a \emph{fixed open service} is a directed walk
inside~$C$ that starts and ends at ports and visits every vertex in
$\mathcal P(C)\cup A$.  Its two endpoint ports are free choices, and repeated
vertices are allowed.  Let $F_k(A)$ be its minimum cost.  Shortcutting repeated
visits in the metric shows that allowing walks does not change this optimum.

For adaptive services we use the following boundary relaxation.  Initially
the policy is in an abstract exterior state.  It must probe every vertex of~$C$
exactly once.  Probes returning an inactive outcome before the first movement
leave the policy in the exterior state.  When the first active vertex~$v$ is probed, the policy
chooses an entry port $a\in\mathcal P(C)$ and is charged $d_C(a,v)$; the
external movement to~$a$ is not charged.  Thereafter the usual adaptive rules
apply with metric~$d_C$.  Once all vertices have been probed, the policy chooses
an exit port $b\in\mathcal P(C)$ and, if its current position is~$x$, is
charged $d_C(x,b)$ for the final internal movement.  Let $G_k$ be the minimum
expected cost of such an
\emph{adaptive open service}, where the expectation is over the independent
midpoint activations in~$C$.  The standard backward-induction argument for the adaptive model shows
that history dependence and private randomization
cannot improve this finite boundary-state problem.  For the single-port block,
\[
    F_0(\emptyset)=G_0=0.
\]

For cost accounting, fix once and for all a shortest generating-graph walk for
each ordered pair of vertices.  Whenever a metric move is made, \emph{expand}
it into this walk.  Passing through a vertex in an expanded walk neither probes
nor serves that vertex.  Every non-root arc occurrence in an expanded walk
belongs to a unique recursive level.

Fix a child block~$C$.  Cut an expanded execution at every crossing of the
boundary of~$C$, including multiple crossings made during a single metric
move.  A \emph{$C$-piece} is a maximal internal portion containing at least
one service event in~$C$; it enters and leaves at ports.  Initial and terminal
port endpoints are treated as boundary crossings.  Internal portions that
only transit through~$C$ may be omitted from the list of pieces, but their
arcs remain included in the total internal cost~$Z$.  Omitting such transit
can only decrease the cost of the local service constructed below.  Probes of
inactive vertices made outside~$C$ cause no movement and do not create pieces.

\begin{lemma}[Ports and interrupted services]\label{lem:block-excursions}
Let $C$ be a copy of~$B_k$.
\begin{enumerate}
    \item Any two ports of~$C$ can be joined inside~$C$ by a directed walk of
    length at most $w_{k+1}=2w_k$ when $k\ge1$; for $k=0$, the block has only
    one port and the corresponding bound is~$0\le w_1$.
    \item Suppose $C$ occurs as a child of a copy of~$B_{k+1}$.  If a
    fixed-realization service is split into $N\ge1$ $C$-pieces, $Z$ is the
    total cost of all arc occurrences internal to~$C$, and
    $A_C\subseteq\mathcal M(C)$ is the active-midpoint set, then
    \begin{equation}\label{eq:post-excursion}
        Z+w_{k+1}(N-1)\ge F_k(A_C).
    \end{equation}
    \item Under the same embedding, suppose an adaptive service is split into
    a random number $N\ge1$ of $C$-pieces, possibly according to observations
    outside~$C$.  Then
    \begin{equation}\label{eq:adapt-excursion}
        \E[Z]+w_{k+1}\E[N-1]\ge G_k.
    \end{equation}
\end{enumerate}
\end{lemma}

\begin{proof}
We first prove the port bound by induction on~$k$.  It is immediate for
$k=0$.  Let $k\ge1$ and let $x,y$ be ports of~$C$.  If they lie in the same
child, induction gives an internal walk of length at most~$w_k$.  If $x$ lies
in the left child and $y$ in the right, the direct level-$k$ arc has length
$w_k$.  In the reverse direction, the two-arc walk through~$m_k$ has length
$2w_k=w_{k+1}$.  This proves the first statement.

For a fixed realization, retain the $N$ pieces in their original order.
Between consecutive pieces, join the preceding exit port to the next entry
port by an internal walk of length at most~$w_{k+1}$.  Together with all
service-containing internal portions, these connectors form a fixed open
service of~$C$.  Its cost is at most $Z+w_{k+1}(N-1)$, proving
\cref{eq:post-excursion}.

For the adaptive statement, let $\xi$ consist of the complete activation
vector outside~$C$, together with any private random seed used by the original
policy.  Sample $\xi$ at time zero.  It is independent of all activations
inside~$C$.  We run a virtual copy of the original execution: probes outside
$C$ receive the corresponding pre-sampled coordinate of~$\xi$ only when that
vertex is virtually probed, and their movements update only the virtual
position in the ambient execution; probes inside~$C$ use the actual local
  outcomes.  Thus every decision depends only on observations revealed inside
  $C$, even though the virtual exterior behavior may depend on them.

Consider two consecutive probes with active outcomes in~$C$.  If they belong to the same
piece, the internal portion of the original expanded movement connects them.
If they belong to different pieces, concatenate three internal walks: the
suffix leading from the first vertex to the preceding exit port, a
port-to-port connector of length at most~$w_{k+1}$, and the prefix from the
next entry port to the second vertex.  The local metric distance between the
two vertices is no larger than this concatenated walk.  The first piece uses
the free-entry convention, and the final internal suffix supplies the terminal
movement to an exit port.  Local probes with inactive outcomes cause no movement.  This
construction also covers arbitrarily many boundary crossings within one
expanded metric move.

For every joint realization, the resulting local strategy has cost at most
$Z+w_{k+1}(N-1)$.  Averaging over the independent seed~$\xi$ and the
activations in~$C$ gives a randomized, possibly history-dependent open service
with expected cost at most the left-hand side of
\cref{eq:adapt-excursion}.  By backward induction for the finite local state
space, a deterministic Markov open service is no more expensive.  Its cost is
at least~$G_k$, proving \cref{eq:adapt-excursion}.
\end{proof}

\subsection{Top-level accounting}\label{sec:recursive-accounting}

At one recursive node, contract the two children to symbols $L$ and $R$ and
write $M$ for the new midpoint.  The level arcs form
\[
    L\longrightarrow R\longrightarrow M\longrightarrow L,
\]
each of unit length after scaling by the common factor~$w$.  In the full
instance add the root symbol~$r$ and the unit arcs
$r\leftrightarrow L$ and $r\leftrightarrow R$.

Let $N_L$ and $N_R$ be the numbers of service pieces in the two children, and
let $H$ be the number of top-level arc \emph{occurrences} in the expanded
execution, counted with multiplicity.  In the closed case, $H$ also counts all
root-arc occurrences, including intermediate passages through~$r$.  Define
\[
    \widehat H:=H-(N_L-1)-(N_R-1).
\]

\begin{lemma}[Top-level accounting]\label{lem:top-accounting}
The following bounds hold.
\begin{enumerate}
    \item For a fixed realization, an open service satisfies
    \begin{equation}\label{eq:open-top-fixed}
        \widehat H\ge1+\1_{\{M\text{ active}\}},
    \end{equation}
    while a service starting and ending at~$r$ satisfies
    \begin{equation}\label{eq:closed-top-fixed}
        \widehat H\ge3+\1_{\{M\text{ active}\}}.
    \end{equation}
    \item For an adaptive service, when $M$ is independently active with
    probability~$1/2$,
    \begin{equation}\label{eq:open-top-adapt}
        \E[\widehat H]\ge2
    \end{equation}
    in the open case, and
    \begin{equation}\label{eq:closed-top-adapt}
        \E[\widehat H]\ge4
    \end{equation}
    in the case starting and ending at~$r$.
\end{enumerate}
\end{lemma}

\begin{proof}
For every piece of $L$ except its first, charge the last top-level arc
occurrence by which the execution enters $L$ before that piece.  Charge pieces
of~$R$ analogously.  These charged occurrences are distinct, one for each
extra piece.  All other top-level arc occurrences are uncharged, so their
number is exactly~$\widehat H$.

In an open fixed-realization service with $M$ inactive, the first entry into
whichever child is not the starting child is uncharged, giving
$\widehat H\ge1$.  If $M$ is active, its service requires another uncharged
entry, giving $\widehat H\ge2$.  In the closed case, the first entries used to
serve $L$ and $R$ and the final return to~$r$ are three distinct uncharged
occurrences.  An active~$M$ requires a fourth.  This proves
\cref{eq:open-top-fixed,eq:closed-top-fixed}.

We next record the equality cases when $M$ is inactive.  Contract every child
service piece to its letter in $\{L,R\}$.  Without the root, the minimum
numbers of top-level arcs between consecutive letters are
\[
    c(L,R)=1,\qquad c(R,L)=2,\qquad
    c(L,L),c(R,R)\ge3.
\]
The piece word contains both letters, and a word with $t$ letters has exactly
$t-2$ charged entries.  If $t=2$, the adjusted costs of the words $LR$ and
$RL$ are respectively~$1$ and~$2$.  If $t\ge3$, the transition table shows
that the sum of transition costs minus $(t-2)$ is at least~$2$: an additional
equal letter costs at least three before its one-unit charge, while any
additional alternation that includes a return $R\to L$ costs two.  Therefore
\begin{equation}\label{eq:open-equality-trace}
    \widehat H=1
    \quad\Longrightarrow\quad
    \text{there is one piece per child and the trace is }L\to R.
\end{equation}

In the closed quotient, entering the first child and returning from the last
each cost one arc.  Intermediate use of the root gives
$c(L,L),c(R,R)\ge2$, while $c(L,R)=1$ and $c(R,L)=2$ remain valid.  For two
pieces, $r\to L\to R\to r$ has adjusted cost~$3$, whereas the reverse order
has cost~$4$.  For $t\ge3$, the transition sum is at least~$t$, so after
adding the two boundary arcs and subtracting $t-2$ charged entries the adjusted
cost is at least~$4$.  Hence
\begin{equation}\label{eq:closed-equality-trace}
    \widehat H=3
    \quad\Longrightarrow\quad
    \text{there is one piece per child and the trace is }
    r\to L\to R\to r.
\end{equation}
These classifications include pure transit, repeated pieces, and intermediate
root passages, since all their top-level arc occurrences were retained in
$H$.

For the adaptive bounds, condition on all activation bits except that of~$M$
and, if present, on the policy's private seed.  Couple the executions with
$M$ inactive and active, and denote their adjusted counts by
$\widehat H_0$ and~$\widehat H_1$.

We will repeatedly use the following elementary counting observation.  If
$t$ specified arc occurrences are unavoidable and at most $s$ of those
occurrences can be charged to extra child pieces, then
$\widehat H\ge t-s$.  Indeed, every other charged occurrence must be one of
the remaining $H-t$ occurrences, so subtracting all charges can cancel at
most $s$ of the specified occurrences.  This observation allows the active
execution to create arbitrarily many additional pieces without weakening the
counts below.

In the open case, the fixed bounds give $\widehat H_0\ge1$ and
$\widehat H_1\ge2$.  If $\widehat H_0\ge2$, their sum is at least~$4$.
Otherwise \cref{eq:open-equality-trace} applies, and the two executions agree
until $M$ is probed.  There are three possible phases.
\begin{itemize}
    \item If $M$ is probed in the exterior state, an active execution needs at
    least the three occurrences $R\to M\to L\to R$.  The entry into~$M$ and
    the first service entries into $L$ and $R$ are all uncharged, so
    $\widehat H_1\ge3$.
    \item If it is probed during the $L$-service, reaching $M$ requires
    $L\to R\to M$, and subsequently serving $R$ requires
    $M\to L\to R$.  Of these four occurrences, only $M\to L$ can be charged,
    namely when an unfinished $L$-service is resumed there.  Hence
    $\widehat H_1\ge4-1=3$; completing~$L$ before the probe cannot reduce the
    count.
    \item If it is probed during or after the $R$-service, the common prefix
    already contains the uncharged first entry $L\to R$.  Serving $M$ and
    reaching a terminal port requires $R\to M\to L$.  All service in~$L$ was
    completed in the common prefix, so $M\to L$ is pure transit or terminal
    movement and cannot be charged.  These three specified occurrences are
    therefore uncharged.  If an unfinished $R$-service is later re-entered,
    its charged entry is an additional occurrence and does not reduce this
    bound.
\end{itemize}
Thus $\widehat H_1\ge3$ whenever $\widehat H_0=1$, and in all cases
\[
    \widehat H_0+\widehat H_1\ge4.
\]
Averaging the equiprobable values of~$M$ proves
\cref{eq:open-top-adapt}.

In the closed case, the fixed bounds give $\widehat H_0\ge3$ and
$\widehat H_1\ge4$.  The desired pair sum is immediate when
$\widehat H_0\ge4$.  If $\widehat H_0=3$, then
\cref{eq:closed-equality-trace} applies.  Again the executions agree until
$M$ is probed.
\begin{itemize}
    \item If $M$ is probed at the initial root, an active execution needs at
    least the five occurrences $r\to R\to M\to L\to R\to r$.  They consist
    of the entry into~$M$, the first service entries into both children, the
    final return, and the necessary transit before~$M$; none is chargeable.
    Thus $\widehat H_1\ge5$.
    \item If it is probed during the $L$-service, the necessary occurrences
    are
    \[
        r\to L\to R\to M\to L\to R\to r.
    \]
    There are six of them, and only $M\to L$ can be charged, when an
    unfinished $L$-service is resumed.  Hence
    $\widehat H_1\ge6-1=5$; if $L$ is already complete, the count is no
    smaller.
    \item If it is probed during or after the $R$-service, the common prefix is
    $r\to L\to R$.  If the $R$-service is complete, the five occurrences in
    $r\to L\to R\to M\to L\to r$ are all uncharged: in particular,
    $M\to L$ is pure transit because the $L$-service is complete.  If the
    $R$-service is unfinished, returning to it and then to~$r$ produces six
    specified occurrences, of which only the re-entry $L\to R$ can be
    charged.  Both subcases give $\widehat H_1\ge5$.
\end{itemize}
There is no case in which $M$ is probed after a final return to~$r$, because
the policy may return only after all vertices have been probed.  Consequently,
$\widehat H_0+\widehat H_1\ge8$, and averaging proves
\cref{eq:closed-top-adapt}.
\end{proof}

\subsection{Gluing relaxed block policies}\label{sec:recursive-gluing}

The open-service definition allows a terminal movement to a chosen port,
whereas the original adaptive model permits movement only after a probe with
an active outcome
or on the final return to the root.  The following observation justifies all
recursive policies used below.

\begin{lemma}[Policy gluing]\label{lem:block-policy-gluing}
Consider an adaptive recursive plan obtained by concatenating adaptive open
services of child blocks with generating-graph walks between their chosen
ports.  The plan may make probes with inactive outcomes before entering a child, and its
entry and exit ports may depend on the observed history.  If every terminal
child movement is included in the stated cost of the plan, then the plan can
be implemented by a policy in the original metric model with no larger cost.
The same conclusion holds for recursively nested concatenations.
\end{lemma}

\begin{proof}
Do not execute a child's terminal movement when its local service ends.
Instead, keep the internal walk from the last active vertex to the chosen exit
port as a pending walk.  Probes with inactive outcomes before the next entry cause no
movement, so the pending walk remains unchanged.  When the next active
vertex~$v$ is probed, concatenate the pending exit walk, the prescribed
parent-level connection to the next entry port, and the next child's internal
entry walk to~$v$.  This is a directed generating-graph walk from the current
active vertex to~$v$; the metric distance charged by the original model is no
larger than its length.  At the end of the full execution, concatenate the
remaining pending walk with the prescribed return to the root.

All ports and connections are chosen from information available when the
corresponding probe returns an active outcome.  The complete arc bundles between child
ports, and between top-level ports and the root, realize every such choice.
The same pending-walk invariant applies at each recursive level, so nested
terminal movements are handled by concatenation rather than by voluntary
repositioning.  Each arc of the recursive plan is charged at most once, which
proves the claim for the resulting adaptive, possibly finite-memory policy.
The standard Markovization argument replaces this policy, without
increasing expected cost, by a deterministic policy of the form $\pi(u,S)$.
\end{proof}

\section{Proof of the a posteriori cost lemma}\label{sec:recursive-posterior-proof}

We evaluate the a posteriori optimum in two steps.  First, a recurrence
expresses the optimal open service of a block through the corresponding
services of its two children.  We then account for the initial departure from
and final return to the root, and average the resulting identity over all
midpoint activations.

\begin{lemma}[Fixed open recurrence]\label{lem:post-rec}
For every $k\ge1$ and every $A\subseteq\mathcal M_k$, let
\[
    A_{\mathrm L}:=A\cap\mathcal M(B_{k-1}^{\mathrm L}),
    \qquad
    A_{\mathrm R}:=A\cap\mathcal M(B_{k-1}^{\mathrm R}).
\]
Then
\begin{equation}\label{eq:post-rec}
    F_k(A)
    =
    F_{k-1}(A_{\mathrm L})+F_{k-1}(A_{\mathrm R})
    +w_k\bigl(1+\1_{\{m_k\in A\}}\bigr).
\end{equation}
\end{lemma}

\begin{proof}
Expand an arbitrary fixed open service.  Let $Z_{\mathrm L}$ and
$Z_{\mathrm R}$ be the total costs of arcs internal to the children,
$N_{\mathrm L}$ and $N_{\mathrm R}$ their numbers of pieces, and $H$ the
number of top-level arc occurrences.  \Cref{lem:block-excursions,lem:top-accounting} give
\begin{align*}
    \cost
    &=Z_{\mathrm L}+Z_{\mathrm R}+w_kH\\
    &\ge F_{k-1}(A_{\mathrm L})+F_{k-1}(A_{\mathrm R})
       +w_k\bigl(H-(N_{\mathrm L}-1)-(N_{\mathrm R}-1)\bigr)\\
    &\ge F_{k-1}(A_{\mathrm L})+F_{k-1}(A_{\mathrm R})
       +w_k\bigl(1+\1_{\{m_k\in A\}}\bigr).
\end{align*}

For the reverse inequality, choose optimal child services together with their
initial and terminal ports.  If $m_k$ is inactive, serve the left child and
use the direct length-$w_k$ port arc to enter and serve the right child.  If
$m_k$ is active, serve the right child, use the two arcs
$R\to m_k\to L$ of total length~$2w_k$, and then serve the left child.  The
complete port bundles allow the required endpoint choices, so both
constructions attain the lower bound.
\end{proof}

\begin{lemma}[Root normalization]\label{lem:recursive-root-normalization}
For every $l\ge1$ and every $A\subseteq\mathcal M_l$,
\begin{equation}\label{eq:recursive-root-normalization}
    \opt(\mathcal P_l\cup A)=2w_l+F_l(A).
\end{equation}
\end{lemma}

\begin{proof}
Expand a tour on $\{r\}\cup\mathcal P_l\cup A$.  Every internal occurrence of
the root in a vertex-to-vertex expanded path has the form
$x\to r\to y$ for ports $x,y\in\mathcal P_l$ and costs~$2w_l$.  By the port
bound in \cref{lem:block-excursions}, it can be replaced by an internal
$x$-to-$y$ walk of length at most $w_{l+1}=2w_l$.  Replace every such
occurrence.  Only the first departure from and the final return to the root
remain.  Each contributes one root arc of length~$w_l$; this remains true
when the first or last served vertex is a midpoint, because its expanded path
enters or leaves $B_l$ through a port.  Removing these two root arcs leaves a
fixed open service, so every tour costs at least $2w_l+F_l(A)$.

Conversely, add a root arc before and after an optimal fixed open service.
The resulting closed walk costs $2w_l+F_l(A)$ and serves every active vertex.
Shortcutting repeated vertex visits in the asymmetric metric produces a tour of
no greater cost.  This proves equality.
\end{proof}

\begin{proof}[Proof of \cref{lem:post-cost}]
Let $\overline F_k:=\E[F_k(A)]$, where the expectation is over the midpoint
activations in~$B_k$.  Since $m_k$ is active with probability~$1/2$,
\[
    \overline F_k
    =
    2\overline F_{k-1}+\frac32w_k,
    \qquad
    \overline F_0=0.
\]
Equivalently, summing over the $2^{k-j}$ level-$j$ midpoints gives
\begin{equation}\label{eq:expected-open-post}
    \overline F_k
    =
    \frac32\sum_{j=1}^k2^{k-j}w_j
    =
    \frac32\sum_{j=1}^k2^{k-1}
    =
    3k2^{k-2}.
\end{equation}
Taking expectations in \cref{eq:recursive-root-normalization} and using
$2w_l=2^l$ yields
\begin{equation}\label{eq:post-value-recursive}
    \optpost(\mathcal{I}_l^{\mathrm{REC}})
    =
    2^l+3l2^{l-2}
    =
    (3l+4)2^{l-2}.
\end{equation}
\end{proof}

\section{Proof of the adaptive cost lemma}\label{sec:recursive-adaptive-proof}

The adaptive analysis follows the same decomposition, but the two executions
obtained from an inactive or active top-level midpoint share a common history
until that midpoint is probed.  The accounting lemmas convert this adaptive
constraint into the open-service recurrence below; we then lift it to the
closed instance.

\begin{lemma}[Adaptive open recurrence]\label{lem:adapt-rec}
For every $k\ge1$,
\begin{equation}\label{eq:adapt-rec}
    G_k=2G_{k-1}+2w_k.
\end{equation}
Consequently, $G_k=k2^k$.
\end{lemma}

\begin{proof}
Consider an arbitrary adaptive open service and use the decomposition from the
proof of \cref{lem:post-rec}.  The adaptive part of
\cref{lem:block-excursions} and the open part of
\cref{lem:top-accounting} give
\begin{align*}
    \E[\cost]
    &=\E[Z_{\mathrm L}+Z_{\mathrm R}+w_kH]\\
    &\ge2G_{k-1}
      +w_k\E\bigl[H-(N_{\mathrm L}-1)-(N_{\mathrm R}-1)\bigr]\\
    &\ge2G_{k-1}+2w_k.
\end{align*}

For the reverse inequality, use optimal adaptive open services in the two
children.  First serve the left child, then cross to and serve the right
child, and finally probe~$m_k$.  If $m_k$ is inactive, terminate at the chosen
right-child exit port.  If it is active, use $R\to m_k\to L$ and terminate at
a left-child port.  \Cref{lem:block-policy-gluing} makes this recursive
description a legal adaptive open service.  Its top-level cost is $w_k$ when
$m_k$ is inactive and $3w_k$ when it is active, hence its expected top-level
cost is~$2w_k$.  Equality follows.  Induction from $G_0=0$ and
$w_k=2^{k-1}$ gives $G_k=k2^k$.
\end{proof}

\begin{proof}[Proof of \cref{lem:adapt-cost}]
Expand an arbitrary adaptive policy for~$\mathcal{I}_l^{\mathrm{REC}}$ and
decompose its work in the two top-level children.
\Cref{lem:block-excursions,lem:top-accounting}, now
using the closed quotient, give
\begin{align*}
    \E[\cost]
    &=\E[Z_{\mathrm L}+Z_{\mathrm R}+w_lH]\\
    &\ge2G_{l-1}
      +w_l\E\bigl[H-(N_{\mathrm L}-1)-(N_{\mathrm R}-1)\bigr]\\
    &\ge2G_{l-1}+4w_l.
\end{align*}

For a matching policy, start at the root, recursively serve the left child
and then the right child, and finally probe~$m_l$.
\Cref{lem:block-policy-gluing} realizes the relaxed child interfaces as legal
metric movements.  If $m_l$ is inactive, the top-level trace is
$r\to L\to R\to r$ and costs~$3w_l$.  If it is active, the trace is
$r\to L\to R\to m_l\to L\to r$ and costs~$5w_l$.  Its expected top-level
cost is therefore~$4w_l$, so the lower bound is attained.  Using
$G_{l-1}=(l-1)2^{l-1}$ and $w_l=2^{l-1}$ gives
\begin{equation}\label{eq:adapt-value-recursive}
    \optadapt(\mathcal{I}_l^{\mathrm{REC}})
    =
    2G_{l-1}+4w_l
    =
    (l+1)2^l.
\end{equation}
\end{proof}

\end{document}
